\evidencelead{E04 · Codex交付 / 外部审核请求}{技术完成之后，我要求独立读取真实成果}{状态汇报是待核查的输入，不是已被用户验收的最终结论。}
\evidenceshot{D04_user_handoff}{我转交的Codex回执}{Codex报告已提交固定GitHub版本，但Process聊天证据尚未补完；工程完成要另行核查。}{15}{119}{38}{C2}
\sourceband{RAW VIDEO 04:51.500 · E04 / USER转交Codex文本；引号内容身份为Codex报告}
\evidenceshot{D04_user_audit}{我的外部审核要求}{我要求GPT从审核包和GitHub独立检查，同步后才给清理建议，并要求本轮先不进入下一阶段。}{9}{119}{29}{C2}
\sourceband{RAW VIDEO 05:05.500 · E04 / USER CONTINUATION · 审核请求与Codex报告在同一用户消息内}
\par\vspace{6mm}\noindent{\sffamily\large\bfseries\color{Ink}GitHub版本身份与技术结论分别验收}\par\vspace{1mm}{\small\color{Muted}先核对来源与版本，再独立复算关键指标；不把未运行的检查写成已核验。}\par\vspace{3mm}
\evidenceshot{D04_gpt_audit}{GPT外审证据}{GPT报告ZIP有832成员、831条文件清单记录通过哈希验证，并直接读取固定GitHub提交。}{25}{119}{42}{C1}
\sourceband{RAW VIDEO 05:08.500 · E04 / GPT AFTER · 对应固定commit: 3ff6353915fc…}
\noindent {\small\color{Muted}第二次审核还包括独立数值复算：在线延迟、车队收益、主动调度成本情景，以及125条进化指标的HV/Spacing。报告明确保留了原始数据全量重训不在本次外部复算范围内。}\par
\vspace{2mm}
\Needspace{37mm}
\evidenceshot{D05_user_status}{我的完成度追问}{我继续询问“除了process report之外”是否完成，以区分技术成果完成、理解验收与教师提交。}{2}{117}{6}{C2}
\sourceband{RAW VIDEO 05:19.500 · E05 / USER · 新互动以用户原生UI开头}
\evidenceshot{D05_gpt_status}{GPT明确区分三类状态}{GPT说明技术和实验成果已完成范围的核验，但Process原生UI、理解验收和教师正式提交不能一起宣布通过。}{13}{118}{21}{C1}
\sourceband{RAW VIDEO 05:20.500 · E05 / GPT AFTER · 对完成度追问的直接回应}
